\documentclass[preprint]{aastex631}

\begin{document}
\title{Methodology and Exploratory Analysis of Stellar Host Parameters and Exoplanet Systems via Principal Component Analysis}
\author{Agastya Gaur}
\noaffiliation

\section{Dataset and Preprocessing}

The data set used for the principal component analysis (PCA) was obtained from the NASA Exoplanet Archive \textit{STELLARHOSTS} table \citep{nasaexoplanetarchivePlanetarySystemsTable2026}\footnote{Accessed on 2026-02-22 at 11:35:01, returning 47,048 rows.}. The set of parameters extracted for this analysis is summarized in \autoref{tab:data_parameters}. The \textit{STELLARHOSTS} table provides stellar parameters for all systems hosting confirmed exoplanets. For systems comprising multiple gravitationally bound stars, the parameters for each component are included.

Because the table includes a separate row for each reference, a host star with multiple published references can appear in multiple rows. To prepare the data for PCA, a data-cleaning procedure was applied to retain only one representative row per host star. The selection procedure was as follows:
\begin{enumerate}
  \item The publication year of each stellar parameter entry was extracted from the \verb|st_refname| field, and a new column, \verb|st_ref_year|, was added to the data set to store this information.

  \item Parameter completeness was computed for each row, defined as the count of non-null values among the following columns: \verb|st_teff|, \verb|st_rad|, \verb|st_mass|, \verb|st_met|, \verb|st_lum|, \verb|st_logg|, \verb|st_age|, \verb|st_dens|, \verb|st_vsin|, \verb|st_rotp|, and \verb|st_radv|.

  \item For each star, the row corresponding to the most recent publication year was retained. In cases of identical or missing publication years, the row with the highest parameter completeness was selected. Remaining ties were resolved by random selection.
  
  \item The ratio used in each rown of \verb|st_met| was defined in \verb|st_metratio|. Based on this, \verb|st_met| was also split into four separate columns: \verb|st_met_FeH|, which contains the [Fe/H] metallicity value, \verb|st_met_MH|, which contains the [M/H] metallicity value, \verb|st_met_NH|, which contains the [N/H] metallicity value, and \verb|st_met_mH|, which contains the [m/H] metallicity value. The \verb|st_met| column was then dropped from the data set.

\end{enumerate}
After applying this filtering procedure, the data set was reduced to 5,008 rows, with each row representing a unique host star and its most recent, most complete set of stellar parameters. The resulting cleaned data set was then used for PCA on the selected stellar parameters.
\begin{table}
  \centering
  \begin{tabular}{ll}
    \hline
    \textbf{Parameter} & \textbf{Description} \\
    \hline
    \verb|sy_name| & System Name \\
    \verb|hostname| & Host Name \\
    \verb|sy_snum| & Number of Stars \\
    \verb|sy_pnum| & Number of Planets \\
    \verb|st_refname| & Stellar Parameter Reference \\
    \verb|st_spectype| & Spectral Type \\
    \verb|st_teff| & Stellar Effective Temperature [K] \\
    \verb|st_rad| & Stellar Radius [Solar Radius] \\
    \verb|st_mass| & Stellar Mass [Solar mass] \\
    \verb|st_met| & Stellar Metallicity [dex] \\
    \verb|st_metratio| & Stellar Metallicity Ratio \\
    \verb|st_lum| & Stellar Luminosity [log(Solar)] \\
    \verb|st_logg| & Stellar Surface Gravity [log10(cm/s$^2$)] \\
    \verb|st_age| & Stellar Age [Gyr] \\
    \verb|st_dens| & Stellar Density [g/cm$^3$] \\
    \verb|st_vsin| & Stellar Rotational Velocity [km/s] \\
    \verb|st_rotp| & Stellar Rotational Period [days] \\
    \verb|st_radv| & Systemic Radial Velocity [km/s] \\
    \hline
  \end{tabular}
  \caption{Data parameters collected from the NASA Exoplanet Archive STELLARHOSTS Table. More information can be found at \url{https://exoplanetarchive.ipac.caltech.edu/docs/API_STELLARHOSTS_columns.html}.}
  \label{tab:data_parameters}
\end{table}

\section{Principal Component Analysis}
Principal component analysis (PCA) is a dimensionality reduction technique that transforms a set of correlated variables into a smaller number of uncorrelated variables called principal components. For the purposes of this project, PCA was applied as an exploratory tool to identify underlying patterns and relationships among the stellar parameters in the dataset. By analyzing the loadings of each principal component, insights can be gained into which parameters contribute most strongly to the variance in the data, potentially revealing clusters of similar host stars or correlations between certain stellar properties.

To perform PCA, a dataset of $n$ items with $p$ features must be represented as a matrix $\mathbf{X} \in \mathbb{R}^{n \times p}$. The columns of $\mathbf{X}$ are then standardized and the covariance matrix $\mathbf{C}$ of the data can be computed:
\begin{equation}
  \mathbf{C} = \frac{1}{n-1} \mathbf{X}^T \mathbf{X}
\end{equation}
The eigenvalues and eigenvectors of $\mathbf{C}$ are then calculated. The eigenvectors define the directions of maximum variance in the data (i.e., the principal component axes), while the corresponding eigenvalues quantify the variance explained by each component. Analyzing the loadings of each component reveals which original features contribute most strongly to the captured variance, thereby offering insight into the underlying structure of the data.

From the \textit{STELLARHOSTS} table, categorical parameters \verb|sy_name|, \verb|hostname|, \verb|st_refname|, \verb|st_spectype|, and \verb|st_metratio| were excluded from the PCA. The issue arose in choosing which of the remaining numerical parameters to include in the analysis, as some parameters had a high proportion of missing values. \autoref{tab:nan_distribution} shows the distribution of missing values across the numerical parameters, sorted by increasing proportion of NaN values. A number of solutions were considered to address the missing data.
\begin{table}
  \centering
  \begin{tabular}{lrr}
    \hline
    \textbf{Parameter} & \textbf{\# NaN} & \textbf{\% NaN} \\
    \hline
    \verb|sy_pnum| & 0 & 0.00\% \\
    \verb|sy_snum| & 0 & 0.00\% \\
    \verb|st_teff| & 999 & 19.95\% \\
    \verb|st_rad| & 1040 & 20.77\% \\
    \verb|st_mass| & 1337 & 26.70\% \\
    \verb|st_met_FeH| & 2395 & 47.82\% \\
    \verb|st_logg| & 2793 & 55.77\% \\
    \verb|st_age| & 2951 & 58.93\% \\
    \verb|st_dens| & 3338 & 66.65\% \\
    \verb|st_lum| & 3942 & 78.71\% \\
    \verb|st_vsin| & 4088 & 81.63\% \\
    \verb|st_radv| & 4590 & 91.65\% \\
    \verb|st_rotp| & 4619 & 92.23\% \\
    \verb|st_met_MH| & 4954 & 98.92\% \\
    \verb|st_met_mH| & 4996 & 99.76\% \\
    \verb|st_met_NH| & 5003 & 99.90\% \\
    \hline
  \end{tabular}
  \caption{Distribution of missing values (NaN counts and percentages) across numerical parameters in the cleaned dataset. Parameters are sorted by increasing proportion of missing values.}
  \label{tab:nan_distribution}
\end{table}

\subsection{Small-sample PCA}
Given the high proportion of missing values in many of the numerical parameters, a small-sample PCA approach was adopted to perform an inital, exploratory analysis. This method involved selecting a subset of the data that contained complete cases for a chosen set of parameters, allowing a trial run of PCA for inital results to be performed without the need for imputation. A completeness threshold $\tau$ was defined, and only parameters with a proportion of missing values less than $\tau$ were included in the PCA. The raw results of all small-sample PCA runs are included in \autoref{sec:appendix_a}.

Small-sample PCA with $\tau = 0.50$ yielded 2,359 rows of data with 5 principal components that together explained 93.21\% of the variance.
\begin{enumerate}
  \item \textbf{Principal Component 1 (24.70\% variance explained)} was strongly influenced by mass (0.72) and radius (0.53), and mildly influenced by temperature (0.34) and the planet count(-0.25). This is a clear main sequence star components, ranging from large, hot, massive stars to smaller, cooler stars. Planet count has in inverse relationship with the size, mass, and temperature. This could be due to observational biases, as larger, hotter stars tend to be more luminous making transit detections and direct imaging difficult. More massive stars also are less influenced by their exoplanets' gravity, making radial velocity detections difficult as well.
  \item \textbf{Principal Component 2 (21.54\% variance explained)} was strongly influenced by temperature (-0.61), metallicity (-0.60), and radius (0.51). This component may show an inverse relationship between stellar host size and temperature/metallicity. This is a strange result, as larger stars are typically hotter and more metal-rich. This can be also attributed to observational bias. Detecting planets around larger, hotter stars is more difficult, so the sample of host stars with measured parameters would not include many of these.
  \item \textbf{Principal Component 3 (17.02\% variance explained)} was strongly influenced by the number of stars (0.81) and the number of planets (-0.52). This component may represent an inverse relationship between systems with more graviationally bound stars and systems with fewer planets, likely due to the gravitational interactions in multi-star systems that can disrupt planet formation and stability.
  \item \textbf{Principal Component 4 (15.83\% variance explained)} was strongly influenced by the number of planets (-0.81) and the number of stars (-0.55). This component captures the difference between sytems with few companion stars and planets and systems with more of both.
  \item \textbf{Principal Component 5 (14.13\% variance explained)} was strongly influenced by metallicity (0.79) and temperature (-0.54). This component clearly represents an inverse relationship between metallicity and temperature. This is a known relationship in stellar astrophysics, recently explored by \citet{martinez-hernandezDESIREDTemperaturemetallicityRelations2026}.
\end{enumerate}

Overall, PC1 and 2 are likely due to observational biases against larger, hotter stars. PC3 and 4 show the variance in exoplanet systems themselves, with PC3 showing the influence of gravitational interactions in multi-star systems and PC4 showing the influence of stellar mass on planet formation. PC5 shows a known relationship between metallicity and temperature. These results suggest that the PCA is capturing meaningful patterns in the data, but also highlights the limitations of the dataset due to missing values and observational biases.

Small-sample PCA with $\tau=0.70$ yielded 583 rows of data with 6 principal components that together explained 95.88\% of the variance. The results were similar to the $\tau=0.50$ case, but with some differences in the loadings and variance explained by each component.
\begin{enumerate}
  \item \textbf{Principal Component 1 (44.53\% variance explained)} was strongly influenced by mass (0.48), radius (0.42), temperature (0.44), log g (-0.46), and mildly influenced by density (-0.37). This is a clear main-sequence star component, ranging from smaller, cooler stars with higher surface gravity and density to larger, hotter stars with lower surface gravity and density. The system parameters contribute weakly, suggesting that most variation in the data is due to intrinsic stellar properties. This component is similar to PC1 in the $\tau=0.50$ case, but with no clear correlation with the number of planets.
  \item \textbf{Principal Component 2 (12.80\% variance explained)} was strongly influenced by age (-0.59), metallicity (-0.53), number of stars (0.49). This component ranges from young, metal-poor stars with more stellar companions to older, metallic stars with fewer companions. This is another astrophysical component, with only weak interactions with the system parameters. This may also be an orthogonal components to the age-metallicity relation, which is a known phenomenon explored by \citet{twarogChemicalEvolutionSolar1980} and \citet{edvardssonChemicalEvolutionGalactic1993}.
  \item \textbf{Principal Component 3 (12.19\% variance explained)} was strongly influenced by the number of planets (-0.74) and the number of stars (0.56). This is a similar component to PC3 in the $\tau=0.50$ case, showing the influence of gravitational interactions in multi-star systems on planet formation and stability.
  \item \textbf{Principal Component 4 (10.44\% variance explained)} was strongly influenced by age (0.67) metallicity (-0.52), and mildly influenced by the number of stars (0.40). This component clearly shows an inverse relationship between age and metallicity and is orthogonal to PC2. This is likely the AMR. 
  \item \textbf{Principal Component 5 (8.70\% variance explained)} was strongly influenced by metallicity (-0.53) number of planets (-0.57), and number of stars (-0.52). This component is similar to PC4 in the $\tau=0.50$ case, simply differentiating between sparse and populated star systems, but also includes shows a strong influence of metallicity proportional to the number of stellar companions and exoplanets.
  \item \textbf{Principal Component 6 (7.12\% variance explained)} was strongly influenced by density (0.60) and radius (0.57), and mildly influenced by log g (-0.31) and temperature (-0.33). This is a weak component that is orthogonal to the main-sequence component (PC1). This is probably present due to differences in stellar structure or evolutionary stage that are not captured by the other parameters.
\end{enumerate}

With more components falling within the threshold, the number sample size decreased significantly, but the results were more astrophysically meaningful. PC1 and PC6 together capture the main-sequence structure of the host stars, while PC2 and PC4 capture the age-metallicity relation. PC3 and PC5 show the influence of stellar companions on planet formation, with PC5 also showing a strong influence of metallicity.

Raising $\tau$ higher than 0.70 greatly diminished the sample size with respect to data dimensionality\footnote{Assuming a well-sampled dataset has $2^D$ samples where $D$ is the number of dimensions.}, and the results were not as meaningful.

Though it contained more samples, $\tau=0.50$ had more components that were likely due to observational biases, while $\tau=0.70$ had fewer samples but more astrophysically meaningful components. This highlights the trade-off between sample size and data quality when performing PCA on datasets with missing values. The small-sample PCA approach provided valuable insights into the underlying structure of the data, but also underscored the limitations of the dataset and the need for careful interpretation of the results.

Both small-sample PCA runs uncovered known astrophysical phenomena and also suggested an inverse relationship between planet count and star count in a system. Though the majority of variance is from intrinsic stellar properties, the system properties also show a non-negligible influence on the variance in the data, suggesting that the dataset contains meaningful patterns related to stellar properties and their influence on exoplanet systems.

This preliminary analysis demonstrated the the likelihood of finding underlying patterns in the stellar host data, but did not directly answer the research question of how stellar host parameters influence exoplanet system properties. To address this, the system parameters can be dropped and small-sample PCA can be performed again to identify the main components of variance in the stellar parameters alone. Then, the system parameters can be projected onto these components to analyze how they relate to the underlying structure of the stellar host properties. This approach would allow for a more direct analysis of the influence of stellar host parameters on exoplanet system properties, while still accounting for the limitations of the dataset.

However, the high proportion of missing values and observational biases limit the conclusions that can be drawn from the analysis. To address this issue, the missing data must be imputed and the full dataset must be artificially extended to perform a more comprehensive PCA on a larger number of parameters. This must be done with observational biases in mind, as the missing data is not random, and the possibility of undetected exoplanets exists for all star systems.

\subsection{Data Imputation}

Missing stellar host data can be attributed to a variety of factors, including observational limitations, selection biases, and the inherent difficulty of measuring certain stellar parameters. For all missing data, true values exist, and their missingness is likely due to their own or other parameters' values.

GSimp, a Gibbs sampler-based imputation method developed by \citet{weiGSimpGibbsSampler2018}, is a powerful tool for imputing missing not at random (MNAR) data in datasets with complex relationships between variables. A modified version of GSimp can be applied to the stellar host dataset to impute missing values while accounting for the underlying structure of the data and the non-random nature of the missingness.

\appendix
\section[Appendix A]{Small-sample PCA Results}
\label{sec:appendix_a}
\subsection[A.1]{PCA with $\tau = 0.50$}
Included parameters: \verb|sy_pnum|, \verb|sy_snum|, \verb|st_teff|, \verb|st_rad|, \verb|st_mass|, and \verb|st_met_FeH|. Data matrix $\mathbf{X} \in \mathbb{R}^{2359 \times 6}$. The raw results are as follows:
\begin{verbatim}
Principal Component 1:
  Eigenvalue: 699.2216
  Proportion of Variance: 24.7005%
  Contribution of Features:
        # of Planets: -0.2487
          # of Stars:  0.1353
               T_eff:  0.3366
              Radius:  0.5344
                Mass:  0.7214
                Fe/H:  0.0242

Principal Component 2:
  Eigenvalue: 609.6782
  Proportion of Variance: 21.5373%
  Contribution of Features:
        # of Planets:  0.0566
          # of Stars:  0.0334
               T_eff: -0.6073
              Radius:  0.5143
                Mass: -0.0643
                Fe/H: -0.5985

Principal Component 3:
  Eigenvalue: 481.8507
  Proportion of Variance: 17.0217%
  Contribution of Features:
        # of Planets: -0.5170
          # of Stars:  0.8144
               T_eff: -0.1505
              Radius: -0.1613
                Mass: -0.1421
                Fe/H:  0.0260

Principal Component 4:
  Eigenvalue: 448.0784
  Proportion of Variance: 15.8287%
  Contribution of Features:
        # of Planets: -0.8079
          # of Stars: -0.5512
               T_eff: -0.0916
              Radius:  0.0524
                Mass: -0.1729
                Fe/H:  0.0493

Principal Component 5:
  Eigenvalue: 399.8554
  Proportion of Variance: 14.1252%
  Contribution of Features:
        # of Planets:  0.1091
          # of Stars:  0.0096
               T_eff: -0.5352
              Radius:  0.2790
                Mass:  0.0524
                Fe/H:  0.7881

Principal Component 6:
  Eigenvalue: 192.1158
  Proportion of Variance: 6.7866%
  Contribution of Features:
        # of Planets: -0.0551
          # of Stars: -0.1161
               T_eff: -0.4477
              Radius: -0.5859
                Mass:  0.6501
                Fe/H: -0.1308
\end{verbatim}

\subsection[A.2]{PCA with $\tau = 0.70$}
Included parameters: \verb|sy_pnum|, \verb|sy_snum|, \verb|st_teff|, \verb|st_rad|, \verb|st_mass|, \verb|st_met_FeH|, \verb|st_logg|, \verb|st_age|, and \verb|st_dens|. Data matrix $\mathbf{X} \in \mathbb{R}^{583 \times 9}$. The raw results are as follows:
\begin{verbatim}
Principal Component 1:
  Eigenvalue: 292.0596
  Proportion of Variance: 44.5298%
  Contribution of Features:
        # of Planets: -0.1036
          # of Stars: -0.0421
               T_eff:  0.4388
              Radius:  0.4173
                Mass:  0.4818
                Fe/H:  0.1299
               log g: -0.4587
                 age: -0.1634
             density: -0.3671
Principal Component 2:
  Eigenvalue: 83.9347
  Proportion of Variance: 12.7974%
  Contribution of Features:
        # of Planets:  0.2363
          # of Stars:  0.4942
               T_eff:  0.1891
              Radius: -0.0552
                Mass:  0.0936
                Fe/H: -0.5313
               log g:  0.1009
                 age: -0.5895
             density:  0.1111
Principal Component 3:
  Eigenvalue: 79.9430
  Proportion of Variance: 12.1888%
  Contribution of Features:
        # of Planets: -0.7382
          # of Stars:  0.5564
               T_eff: -0.0709
              Radius: -0.0334
                Mass:  0.0087
                Fe/H:  0.3287
               log g:  0.0594
                 age: -0.1103
             density:  0.1244
Principal Component 4:
  Eigenvalue: 68.4983
  Proportion of Variance: 10.4438%
  Contribution of Features:
        # of Planets: -0.0742
          # of Stars:  0.4042
               T_eff: -0.0301
              Radius:  0.2025
                Mass: -0.0860
                Fe/H: -0.5162
               log g: -0.2252
                 age:  0.6667
             density: -0.1420
Principal Component 5:
  Eigenvalue: 57.0557
  Proportion of Variance: 8.6992%
  Contribution of Features:
        # of Planets: -0.5663
          # of Stars: -0.5204
               T_eff:  0.0104
              Radius:  0.1387
                Mass: -0.0239
                Fe/H: -0.5281
               log g: -0.0412
                 age: -0.1531
             density:  0.2909
Principal Component 6:
  Eigenvalue: 46.6860
  Proportion of Variance: 7.1181%
  Contribution of Features:
        # of Planets:  0.2415
          # of Stars:  0.0993
               T_eff: -0.3310
              Radius:  0.5747
                Mass: -0.0504
                Fe/H:  0.1665
               log g: -0.3135
                 age: -0.0790
             density:  0.5977
Principal Component 7:
  Eigenvalue: 22.6299
  Proportion of Variance: 3.4503%
  Contribution of Features:
        # of Planets: -0.0622
          # of Stars: -0.0159
               T_eff: -0.5741
              Radius:  0.2306
                Mass: -0.3287
                Fe/H: -0.0789
               log g: -0.0805
                 age: -0.3531
             density: -0.6064
Principal Component 8:
  Eigenvalue: 3.2212
  Proportion of Variance: 0.4911%
  Contribution of Features:
        # of Planets: -0.0038
          # of Stars:  0.0033
               T_eff: -0.0534
              Radius: -0.5767
                Mass: -0.1679
                Fe/H: -0.0050
               log g: -0.7863
                 age: -0.1023
             density:  0.0872
Principal Component 9:
  Eigenvalue: 1.8466
  Proportion of Variance: 0.2815%
  Contribution of Features:
        # of Planets:  0.0069
          # of Stars: -0.0105
               T_eff: -0.5690
              Radius: -0.2134
                Mass:  0.7825
                Fe/H: -0.1165
               log g:  0.0210
                 age:  0.0656
             density:  0.0071
\end{verbatim}

\bibliography{references}

\end{document}